\documentclass[10pt,a4paper]{article}

% --- Paketler ---
\usepackage[utf8]{inputenc}
\usepackage[english]{babel}
\usepackage{amsmath, amssymb, amsthm}
\usepackage{geometry}
\geometry{margin=1in}
\usepackage{graphicx}
\usepackage{setspace}
\usepackage{listings}
\usepackage{booktabs}
\usepackage{longtable}
\usepackage[colorlinks=true, linkcolor=black, urlcolor=blue]{hyperref}
\usepackage{xcolor}

% --- Section numbering: 4.1, 4.2, ... 4.8 ---
\renewcommand{\thesection}{4.\arabic{section}}

\lstdefinelanguage{WorkflowScript}{
    keywords={task, pipeline, stage, on_failure, fun, var, run, return,
              if, else, while, int, float, bool, void, timeout, retries,
              parallel, true, false},
    sensitive=true,
    morecomment=[l]{//},
}

\lstset{
    language=WorkflowScript,
    basicstyle=\ttfamily\small,
    keywordstyle=\color{blue},
    commentstyle=\color{gray}\itshape,
    breaklines=true,
    frame=single,
    showstringspaces=false,
    captionpos=b,
}
\onehalfspacing

% --- Teorem Ortamları ---
\newtheorem{theorem}{Theorem}[section]
\newtheorem{definition}{Definition}[section]
\newtheorem{example}{Example}[section]

% --- Kapak Bilgileri ---
\title{\textbf{WorkflowScript: D1 -- Design Specification}\\[0.5em]
       \large CSE341 --- Concepts of Programming Languages\\[0.3em]
       \normalsize Part 1 \& 2 Submission --- 22 May 2026}
\author{Ömer Faruk Koç \\ \small 210104004061 \\ \small Gebze Technical University}
\date{May, 2026}

\begin{document}

\maketitle
\newpage

\tableofcontents
\newpage

% -----------------------------------------------------------------------
\section{Language Overview}
% -----------------------------------------------------------------------

\subsection{Abstract}

WorkflowScript is a domain-specific language for defining, composing,
and executing task-based workflow pipelines. Its target users are backend
developers who need to model CI/CD pipelines, job scheduling, and
multi-step automation without embedding that logic in a general-purpose
language.
The language prioritizes reliability through strong static typing, static scoping, and mandatory task-field declarations, while maintaining writability with concise syntax for common workflow patterns.

YAML-based pipeline definitions (GitHub Actions, GitLab CI) are untyped
and unstructured: a misspelled field or a wrong value type produces a
runtime failure, not a definition-time error. General-purpose languages
add boilerplate and require recompilation for every change.
WorkflowScript occupies the gap: it has enough of a programming model
(functions, variables, control flow) to be expressive, while its type
system and fixed task structure catch errors before any execution begins.

\subsection{Design Priorities}

WorkflowScript prioritizes the following Sebesta criteria:

\begin{itemize}
    \item \textbf{Reliability} --- strong static typing, static scoping,
          and mandatory task fields (\texttt{timeout}, \texttt{retries},
          \texttt{parallel}) prevent undefined behavior at pipeline
          definition time. A program that passes the parser is guaranteed
          to have well-formed task declarations; a program that passes the
          type checker is guaranteed to reference only declared tasks and
          well-typed variables.

    \item \textbf{Writability} --- concise syntax for the common case:
          a three-field \texttt{task} declaration, a one-keyword
          \texttt{stage} reference, and a single \texttt{run} statement
          to trigger a pipeline or task.

    \item \textbf{Readability} --- structure mirrors the mental model of
          a workflow: tasks declared first, pipeline composition second,
          imperative control logic last. The keyword
          \texttt{on\_failure} reads like plain English.
\end{itemize}

Knowingly sacrificed: \textbf{cost} (no string type, no arrays, no
dynamic dispatch, no generics). These features are not needed in the
pipeline domain and would complicate the type checker and interpreter
without enabling any new workflow patterns.

\subsection{Sample Program}

\begin{lstlisting}[caption={Sample WorkflowScript program}]
task build {
    timeout:  30
    retries:  0
    parallel: false
}

task test {
    timeout:  60
    retries:  2
    parallel: true
}
\end{lstlisting}

% -----------------------------------------------------------------------
\section{Lexical Structure}
% -----------------------------------------------------------------------

\subsection{Token Categories}

\begin{table}[h!]
\centering
\begin{tabular}{@{}lll@{}}
\toprule
\textbf{Category} & \textbf{Pattern} & \textbf{Examples} \\
\midrule
IDENT      & \texttt{[a-zA-Z\_][a-zA-Z0-9\_]*}     & \texttt{build}, \texttt{my\_task}, \texttt{ci} \\
INT\_LIT   & \texttt{[0-9]+}                         & \texttt{0}, \texttt{30}, \texttt{120} \\
FLOAT\_LIT & \texttt{[0-9]+ "." [0-9]+}             & \texttt{0.5}, \texttt{3.14}, \texttt{60.0} \\
BOOL\_LIT  & \texttt{"true" | "false"}               & \texttt{true}, \texttt{false} \\
KEYWORD    & Reserved identifiers (see below)        & \texttt{task}, \texttt{if}, \texttt{while} \\
OPERATOR   & Fixed character sequences               & \texttt{+}, \texttt{==}, \texttt{->} \\
SEPARATOR  & Fixed single characters                 & \texttt{(}, \texttt{\{}, \texttt{:} \\
COMMENT    & \texttt{"//" .* "{\textbackslash}n"}    & \texttt{// build the service} \\
\bottomrule
\end{tabular}
\caption{Token categories in WorkflowScript}
\end{table}

\subsection{Keywords}

The following identifiers are reserved and cannot be used as
user-defined names:

\begin{lstlisting}[language={}, frame=none, basicstyle=\ttfamily\small]
task       pipeline   stage      on_failure
fun        var        run        return
if         else       while
int        float      bool       void
timeout    retries    parallel
true       false
\end{lstlisting}

\texttt{true} and \texttt{false} are lexed as BOOL\_LIT tokens rather
than generic keywords, which simplifies the parser: they can appear
anywhere a literal is expected without a separate production rule.
\texttt{task} is reused as both a declaration keyword and a type name in
function parameter lists. \texttt{timeout}, \texttt{retries}, and
\texttt{parallel} are reserved so they cannot shadow task field names in
expressions like \texttt{t.timeout}.

\subsection{Operators}

\begin{lstlisting}[language={}, frame=none, basicstyle=\ttfamily\small]
+    -    *    /
==   !=   <    <=   >    >=
&&   ||   !
=    ->   .
\end{lstlisting}

\subsection{Separators}

\begin{lstlisting}[language={}, frame=none, basicstyle=\ttfamily\small]
(    )    {    }    :    ,
\end{lstlisting}

\subsection{Comments}

Line comments begin with \texttt{//} and extend to the end of the line.
There are no block comments. The lexer discards comments entirely; they
produce no tokens. Whitespace (spaces, tabs, carriage returns, newlines)
is also ignored between tokens.

% -----------------------------------------------------------------------
\section{Syntax}
% -----------------------------------------------------------------------

\subsection{Complete Grammar}

The complete grammar in EBNF is given below.

\begin{lstlisting}[language={}, caption={WorkflowScript EBNF grammar}]
<program>        ::= { <top_decl> }

<top_decl>       ::= <task_decl>
                   | <pipeline_decl>
                   | <fun_decl>
                   | <stmt>

<task_decl>      ::= "task" IDENT "{"
                         "timeout"  ":" INT_LIT
                         "retries"  ":" INT_LIT
                         "parallel" ":" BOOL_LIT
                     "}"

<pipeline_decl>  ::= "pipeline" IDENT "{" <pipeline_body> "}"
<pipeline_body>  ::= <stage_stmt> { <stage_stmt> } [ <on_failure> ]
<stage_stmt>     ::= "stage" IDENT
<on_failure>     ::= "on_failure" "{" { <stmt> } "}"

<fun_decl>       ::= "fun" IDENT "(" [ <params> ] ")" "->" <type> <block>
<params>         ::= <param> { "," <param> }
<param>          ::= IDENT ":" <type>

<type>           ::= "int" | "float" | "bool" | "void" | "task"

<block>          ::= "{" { <stmt> } "}"

<stmt>           ::= <var_decl>
                   | <assign_stmt>
                   | <if_stmt>
                   | <while_stmt>
                   | <run_stmt>
                   | <return_stmt>
                   | <expr_stmt>

<var_decl>       ::= "var" IDENT ":" <type> "=" <expr>
<assign_stmt>    ::= IDENT "=" <expr>
<if_stmt>        ::= "if" <expr> <block> [ "else" <block> ]
<while_stmt>     ::= "while" <expr> <block>
<run_stmt>       ::= "run" IDENT
<return_stmt>    ::= "return" <expr>
<expr_stmt>      ::= <call_expr>

<expr>           ::= <or_expr>
<or_expr>        ::= <and_expr> { "||" <and_expr> }
<and_expr>       ::= <eq_expr> { "&&" <eq_expr> }
<eq_expr>        ::= <cmp_expr> { ( "==" | "!=" ) <cmp_expr> }
<cmp_expr>       ::= <add_expr> { ( "<" | "<=" | ">" | ">=" ) <add_expr> }
<add_expr>       ::= <mul_expr> { ( "+" | "-" ) <mul_expr> }
<mul_expr>       ::= <unary_expr> { ( "*" | "/" ) <unary_expr> }
<unary_expr>     ::= ( "-" | "!" ) <unary_expr> | <primary>
<primary>        ::= INT_LIT
                   | FLOAT_LIT
                   | BOOL_LIT
                   | <call_expr>
                   | <field_expr>
                   | IDENT
                   | "(" <expr> ")"

<call_expr>      ::= IDENT "(" [ <args> ] ")"
<args>           ::= <expr> { "," <expr> }
<field_expr>     ::= IDENT "." <field_name>
<field_name>     ::= IDENT | "timeout" | "retries" | "parallel"
\end{lstlisting}

\subsection{Ambiguity Notes}

\subsubsection{Dangling Else}

The \texttt{<if\_stmt>} rule with an optional \texttt{else} is
potentially ambiguous when \texttt{if} statements are nested.
WorkflowScript resolves this with the nearest-\texttt{if} rule: an
\texttt{else} clause always binds to the innermost unmatched \texttt{if}.
This is the same convention used by C, Go, and Java. The
recursive-descent parser implements it naturally: after parsing the
\texttt{then} block, it looks for \texttt{else} before returning, so the
\texttt{else} is consumed by the current (innermost) \texttt{if}
invocation.

\subsubsection{Operator Precedence}

Precedence is encoded directly into the grammar stratification rather
than via a separate table. From lowest to highest binding strength:

\begin{table}[h!]
\centering
\begin{tabular}{@{}ll@{}}
\toprule
\textbf{Level} & \textbf{Operators} \\
\midrule
Lowest  & \texttt{||} \\
        & \texttt{\&\&} \\
        & \texttt{==} \quad \texttt{!=} \\
        & \texttt{<} \quad \texttt{<=} \quad \texttt{>} \quad \texttt{>=} \\
        & \texttt{+} \quad \texttt{-} \\
        & \texttt{*} \quad \texttt{/} \\
Highest & unary \texttt{-} \quad \texttt{!} \\
\bottomrule
\end{tabular}
\caption{Operator precedence levels (lowest to highest)}
\end{table}

Because each grammar level calls the level above it, there is no
ambiguity about how an expression like
\texttt{a + b * c == d || e} parses --- the structure is determined
by the grammar alone.

\subsubsection{Field Access on Task-Field Keywords}

The identifiers \texttt{timeout}, \texttt{retries}, and \texttt{parallel}
are lexed as keyword tokens, not as IDENT. A naive grammar would make
\texttt{t.timeout} unparseable because \texttt{.} would be followed by
a keyword rather than an IDENT. The \texttt{<field\_name>} non-terminal
explicitly allows all three task-field keywords on the right-hand side
of \texttt{.}, resolving this.

\subsubsection{Expression Statement Restriction}

Only call expressions are permitted as standalone statements
(\texttt{<expr\_stmt> ::= <call\_expr>}). A bare variable reference
used as a statement (e.g., writing \texttt{x} on its own line) is
rejected with a parser error. This eliminates the syntactic ambiguity
between a no-op expression statement and a misspelled assignment, and
prevents silent no-op bugs.

\subsubsection{Assignment vs.\ Expression}

Assignment is a statement (\texttt{<assign\_stmt>}), not an expression.
The parser distinguishes \texttt{IDENT "="} (assignment) from
\texttt{IDENT "("} (call) by one token of lookahead. This prevents the
classic \texttt{if x = 5} bug and removes the need for a separate
assignment-expression production.

% -----------------------------------------------------------------------
\section{Semantics}
% -----------------------------------------------------------------------

Operational semantics (Sebesta \S3.5) is used to describe the meaning of WorkflowScript
constructs. The approach follows the method Sebesta describes: each language construct is
translated into an equivalent sequence of statements in a simpler \textbf{intermediate
language} whose meaning is self-evident. The human reader acts as the virtual machine that
executes these intermediate statements.

\subsection{Intermediate Language}

The intermediate language used below consists of the following statement forms. All statement
forms are drawn from the vocabulary Sebesta \S3.5.1 introduces:

\begin{table}[h!]
\centering
\begin{tabular}{@{}ll@{}}
\toprule
\textbf{Intermediate Statement} & \textbf{Meaning} \\
\midrule
\texttt{ident = expr}            & Assign the value of \texttt{expr} to variable \texttt{ident} \\
\texttt{goto label}              & Unconditional jump to \texttt{label} \\
\texttt{if not(expr) goto label} & Jump to \texttt{label} if \texttt{expr} is false \\
\texttt{output(str)}             & Print \texttt{str} to standard output (side effect) \\
\texttt{error(str)}              & Halt with error message \texttt{str} \\
\bottomrule
\end{tabular}
\caption{Intermediate language statement forms}
\end{table}

Labels are written as \texttt{label:} at the start of a line. \texttt{T} and \texttt{P}
denote the global, immutable task store and pipeline store respectively, both populated
before any statement executes.

\subsection{While Statement}

The \texttt{while} statement is the only looping construct in WorkflowScript. Its meaning is
given by translating it into the intermediate language. The translation has two parts: a
conditional check at the top and an unconditional back-edge at the bottom, matching the
pattern Sebesta \S3.5.1 uses for the C \texttt{for} construct.

\begin{lstlisting}[language={}, caption={While statement translation to intermediate language}]
WorkflowScript                   Intermediate Language
--------------------------       ----------------------------------------
while cond {                     loop: if not(cond) goto end
    s1                                 s1
    s2                                 s2
    ...                                ...
    sn                                 sn
}                                      goto loop
                                 end:
\end{lstlisting}

\textbf{Reading the translation.}
On each iteration the interpreter evaluates \texttt{cond}. If it is false, control jumps to
\texttt{end} and the loop is skipped (or exits). If it is true, the body statements
$s_1 \ldots s_n$ execute in order and control returns to \texttt{loop} for the next test.

\textbf{Non-termination.}
If \texttt{cond} never becomes false, the \texttt{goto loop} back-edge fires indefinitely
and the program diverges. WorkflowScript provides no \texttt{break} or \texttt{continue};
termination is the programmer's responsibility. The type checker guarantees \texttt{cond}
is of type \texttt{bool} but cannot determine whether the loop terminates.

\textbf{Nested blocks.}
Each statement $s_i$ in the body may itself be a block; if so, it is expanded recursively
by the same translation rules. An empty body translates to only the \texttt{goto loop}
back-edge, producing a spin-loop that depends on a side-effecting condition.

\subsection{Run Statement (Domain-Specific Construct)}

\texttt{run} is WorkflowScript's primary domain-specific statement. It triggers either a
single named task or an entire named pipeline. There are two cases depending on whether
\texttt{name} resolves to a task or a pipeline in the global stores.

\textbf{Case 1 --- Running a task}
(\texttt{name} $\in \mathrm{dom}(\mathtt{T})$):

\begin{lstlisting}[language={}, caption={Run task translation}]
WorkflowScript                   Intermediate Language
--------------------------       ----------------------------------------
run name                         output("[run] name")
\end{lstlisting}

A single output line is produced. The environment (variable bindings) is not modified;
\texttt{run} is a pure output action.

\textbf{Case 2 --- Running a pipeline}
(\texttt{name} $\in \mathrm{dom}(\mathtt{P})$, stages $= [t_1, t_2, \ldots, t_n]$):

\begin{lstlisting}[language={}, caption={Run pipeline translation}]
WorkflowScript                   Intermediate Language
--------------------------       ----------------------------------------
run name                         output("[run] t1")
                                 output("[run] t2")
                                 ...
                                 output("[run] tn")
\end{lstlisting}

The pipeline's stage list is unrolled into one \texttt{output} call per stage in declaration
order. There is no conditional branching; every stage always executes.

\textbf{Case 3 --- Undefined name}
(\texttt{name} $\notin \mathrm{dom}(\mathtt{T})$ and \texttt{name} $\notin \mathrm{dom}(\mathtt{P})$):

\begin{lstlisting}[language={}, caption={Run undefined name translation}]
WorkflowScript                   Intermediate Language
--------------------------       ----------------------------------------
run name                         error("undefined task or pipeline 'name'")
\end{lstlisting}

This case is unreachable in well-typed programs: the static type checker verifies that every
\texttt{run} target is declared before execution begins, so \texttt{error(...)} is included
here only for completeness.

\textbf{Key observations.}
\begin{enumerate}
    \item \texttt{run} does \textbf{not} modify variable bindings. Task parameters
          (\texttt{timeout}, \texttt{retries}, \texttt{parallel}) are fixed at declaration
          time; the intermediate translation contains no assignment to any program variable.
    \item Case~1 vs.\ Case~2 is resolved by checking $\mathrm{dom}(\mathtt{T})$ first.
          A name cannot simultaneously be a task and a pipeline because both share the same
          global namespace and the parser rejects duplicate declarations.
    \item The \texttt{on\_failure} block of a \texttt{pipeline} declaration is statically
          type-checked but is never entered during normal execution: since \texttt{run} is
          simulated (output only), no failure signal is ever raised. It would be triggered
          by an extension that models real process failure.
\end{enumerate}

% -----------------------------------------------------------------------
\section{Type System}
% -----------------------------------------------------------------------

\subsection{Primitive Types and Value Sets}

WorkflowScript has three primitive types:

\begin{table}[h!]
\centering
\begin{tabular}{@{}lll@{}}
\toprule
\textbf{Type} & \textbf{Go Runtime Representation} & \textbf{Value Set} \\
\midrule
\texttt{int}   & \texttt{int64}   & 64-bit signed integers stored in two's-complement; range $-2^{63}$ to $2^{63}-1$ \\
\texttt{float} & \texttt{float64} & IEEE~754 double precision; $\approx \pm 1.8 \times 10^{308}$, 15--17 sig.\ digits \\
\texttt{bool}  & \texttt{bool}    & \{\texttt{true}, \texttt{false}\} \\
\bottomrule
\end{tabular}
\caption{Primitive types in WorkflowScript}
\end{table}

\texttt{int} covers all realistic timeout values (seconds) and retry counts. \texttt{float}
is available for computations that mix integer task parameters into fractional results.
\texttt{bool} corresponds to the \texttt{parallel} task field and to all comparison results.

There is no \texttt{string} type. Pipeline stage names are identifiers in the source text, not
runtime data; no string manipulation is needed in the workflow domain.

\subsection{Strong or Weak Typing?}

WorkflowScript is \textbf{strongly typed} (Sebesta \S6.12 definition: a language is strongly
typed if type errors are always detected). Every variable carries a declared type annotation
that does not change across assignments. The type checker enforces the following rules before
any execution begins:

\begin{itemize}
    \item A variable declared as \texttt{int} cannot be assigned a \texttt{bool} value.
    \item Arithmetic operators (\texttt{+}, \texttt{-}, \texttt{*}, \texttt{/}) require both
          operands to be numeric (\texttt{int} or \texttt{float}); applying them to
          \texttt{bool} is a compile-time error.
    \item Boolean operators (\texttt{\&\&}, \texttt{||}, \texttt{!}) require \texttt{bool}
          operands.
    \item A function declared \texttt{-> int} must return an expression of type \texttt{int}
          (or \texttt{float} --- see \S\ref{subsec:coercion}); returning \texttt{bool} is a
          compile-time error.
    \item \texttt{void} cannot be used as a value: \texttt{var x: void} is rejected, and the
          result of a \texttt{void} function cannot appear in an expression.
\end{itemize}

No type information is deferred to runtime: every name's type is known statically, and every
expression's type is determined from its sub-expressions without executing the program.

\textbf{Caveat --- coercion and strong typing.}
Sebesta \S6.14 notes that implicit coercion is the primary way languages weaken the
strong-typing guarantee: if an \texttt{int} is silently accepted where a \texttt{float} is
expected, a programmer who writes the wrong variable by mistake will not receive a type error.
WorkflowScript's single coercion rule (int$\to$float widening, \S\ref{subsec:coercion})
introduces this theoretical weakening. In practice the risk is minimal: widening only applies
in one direction and only between numeric types, and the domain is narrow enough that
accidentally passing an integer where a float is expected is unlikely to cause a semantic
error. Nevertheless, the language is best described as \textit{effectively strongly typed}
rather than \textit{absolutely strongly typed} in Sebesta's strictest sense.

\subsection{Implicit Type Coercion}
\label{subsec:coercion}

There is exactly one coercion rule:

\begin{quote}
\textbf{int-to-float widening:} an expression of type \texttt{int} is implicitly accepted
wherever a \texttt{float} is expected, in both variable declarations and function arguments.
\end{quote}

\begin{lstlisting}[language={}, caption={Type compatibility function}]
typesCompatible(int, float) = true
typesCompatible(float, int) = false   -- narrowing is forbidden
typesCompatible(T, T)       = true    -- reflexive for all other types
typesCompatible(T, U)       = false   -- otherwise
\end{lstlisting}

This coercion is applied in:

\begin{itemize}
    \item \textbf{Variable declarations:} \texttt{var x: float = 1 + 2} is valid; the integer
          result \texttt{3} widens to \texttt{3.0}.
    \item \textbf{Function arguments:} a parameter of type \texttt{float} may receive an
          \texttt{int} argument.
    \item \textbf{Binary arithmetic:} if one operand is \texttt{float} and the other is
          \texttt{int}, the \texttt{int} is promoted to \texttt{float} and the result is
          \texttt{float}.
\end{itemize}

Narrowing (\texttt{float}$\to$\texttt{int}) is deliberately absent. Truncation loses
information silently and conflicts with WorkflowScript's reliability goal. Any narrowing
conversion must be written explicitly by the programmer (not supported in the current version,
consistent with the language's minimal scope).

\subsection{Type Equivalence}

WorkflowScript uses \textbf{name equivalence} for the \texttt{task} structured type
(Sebesta \S6.14).

There is exactly one type named \texttt{task}. Every task declaration --- regardless of the
values assigned to \texttt{timeout}, \texttt{retries}, and \texttt{parallel} --- belongs to
this single type. Two tasks with identical field values are therefore type-compatible; a
function accepting \texttt{task} may receive any declared task as an argument.

Name equivalence is the right choice for two reasons.

\textbf{Primary reason --- forward extensibility and semantic safety.}
Sebesta \S6.14 illustrates the hazard of structural equivalence with the classic
Celsius/Fahrenheit example: two record types with an identical single \texttt{float} field
are structurally equivalent, yet treating a Celsius value as Fahrenheit is a semantic error
that structural equivalence silently permits. The same hazard applies here. If a second
structured type --- say, \texttt{pipeline} --- were later added to WorkflowScript, it would
likely carry the same fields (\texttt{timeout: int}, \texttt{retries: int},
\texttt{parallel: bool}). Under structural equivalence a function that expected a
\texttt{task} would silently accept a \texttt{pipeline}, conflating two fundamentally
different domain concepts. Name equivalence prevents this: the type checker would reject the
mismatch regardless of field layout. The principle is that structurally similar types can be
semantically incompatible, and name equivalence is the mechanism that enforces this separation.

\textbf{Secondary reason --- conceptual fit.}
Tasks are domain roles identified by name (\texttt{build}, \texttt{deploy}, \texttt{test}),
not by their parameter values. Two tasks with identical timeouts and retry counts are still
different units of work. Name equivalence reflects this: the type system recognises only the
declared type name \texttt{task}, not the field values, consistent with the domain model.

\subsection{Structured Type Design Decisions (Sebesta \S6.5--\S6.7)}

The \texttt{task} type is WorkflowScript's single structured type. It is a \textbf{record}
(Sebesta \S6.7) with three fixed fields:

\begin{table}[h!]
\centering
\begin{tabular}{@{}lll@{}}
\toprule
\textbf{Field} & \textbf{Type} & \textbf{Meaning} \\
\midrule
\texttt{timeout}  & \texttt{int}  & Maximum seconds the task may run \\
\texttt{retries}  & \texttt{int}  & Number of retry attempts on failure \\
\texttt{parallel} & \texttt{bool} & Whether the task may run concurrently \\
\bottomrule
\end{tabular}
\caption{Fields of the \texttt{task} structured type}
\end{table}

Sebesta \S6.7 identifies two primary design issues for record types; a third concern
(implementation) is addressed in the same section. Each is resolved below for \texttt{task}.

\textbf{Design issue 1 --- Field reference syntax.}
Sebesta \S6.7 contrasts the dot notation used in most languages (\texttt{record.field}) with
COBOL's reversed \texttt{OF} syntax (\texttt{field OF record}). WorkflowScript uses dot
notation exclusively: \texttt{t.timeout}, \texttt{t.retries}, \texttt{t.parallel}. Dot
notation was chosen because it is the near-universal convention in languages that
WorkflowScript programmers are likely to know (Go, Python, Java, JavaScript). The three field
names (\texttt{timeout}, \texttt{retries}, \texttt{parallel}) are reserved keywords so that
the lexer can tokenise \texttt{t.timeout} unambiguously.

\textbf{Design issue 2 --- Elliptical references.}
In languages with nested records, Sebesta \S6.7 notes that some designs allow the programmer
to omit intermediate record names when the field name is unambiguous in scope (an
\textit{elliptical reference}). WorkflowScript deliberately \textbf{disallows} elliptical
references. Every field access must be written as \texttt{variable.field}
(e.g., \texttt{t.timeout}), never as a bare \texttt{timeout}. This decision keeps name
resolution simple: the parser never needs to search the environment for a record that owns a
given field name. It also prevents shadowing bugs where a local variable named \texttt{timeout}
might be confused with a task field of the same name.

\textbf{Implementation --- compile-time descriptor, no run-time descriptor.}
Sebesta \S6.7 explains that record field offsets are fixed at compile time and can be stored
in a \textit{type descriptor}. No descriptor is needed at run time because field names are
never looked up dynamically. WorkflowScript follows this model exactly: the type checker
resolves \texttt{t.timeout} to the \texttt{timeout} field of the \texttt{task} type at
compile time and rejects any unknown field name (e.g., \texttt{t.cost}) as a type error. The
interpreter then accesses \texttt{TaskValue} fields directly by name in Go struct literals ---
there is no hash-map lookup or dynamic dispatch at run time. Field access is O(1) and cannot
fail after the type-checking pass.

% -----------------------------------------------------------------------
\section{Names, Binding, Scope, and Lifetime}
% -----------------------------------------------------------------------

\subsection{Legal Identifiers}

A name must match the regular expression
\texttt{[a-zA-Z\_][a-zA-Z0-9\_]*}: it begins with a letter or
underscore and continues with any mix of letters, digits, and
underscores. There is no length limit. Names are case-sensitive:
\texttt{Build} and \texttt{build} are distinct identifiers. Any name
that matches a keyword listed in Section~4.2 is reserved and cannot be
used as a user-defined identifier.

\subsection{Bindings: Compile Time vs.\ Run Time}

\begin{table}[h!]
\centering
\begin{tabular}{@{}p{7.5cm}p{6.5cm}@{}}
\toprule
\textbf{Binding} & \textbf{Time} \\
\midrule
Variable type (declared in \texttt{var})
    & Compile time --- explicit type annotation \\
Function parameter types and return type
    & Compile time --- explicit annotations \\
Task field types (\texttt{timeout}: int, \texttt{retries}: int, \texttt{parallel}: bool)
    & Compile time --- fixed by language definition \\
Stage-to-task resolution in pipelines
    & Compile time --- stage IDENT must name a declared task \\
Variable value
    & Run time \\
Function body execution
    & Run time \\
Pipeline execution (on \texttt{run})
    & Run time \\
\bottomrule
\end{tabular}
\caption{Binding times in WorkflowScript}
\end{table}

The early binding of types and stage references is the primary mechanism
by which WorkflowScript achieves reliability. A program that passes the
type checker is guaranteed to have no unresolved task references at
runtime, and no type mismatches in arithmetic or comparisons.

\subsection{Scoping}

WorkflowScript uses \textbf{static (lexical) scoping}. The scope of a
name is determined by its textual position in the source. A name is
visible from its declaration point to the end of the enclosing block.
Inner blocks may shadow outer names.

There are three kinds of scope:

\begin{enumerate}
    \item \textbf{Program scope} --- the outermost scope. Task
          declarations, pipeline declarations, top-level \texttt{var}
          declarations, and top-level statements all live here. Every
          name declared at program scope is visible throughout the entire
          program after its declaration point.

    \item \textbf{Function scope} --- each \texttt{fun} body opens a new
          scope. Parameters are bound at the function entry point and are
          visible throughout the function body. \texttt{var} declarations
          inside the body extend their scope to the closing \texttt{\}}
          of the function.

    \item \textbf{Block scope} --- \texttt{if} bodies, \texttt{while}
          bodies, and \texttt{on\_failure} bodies each open a nested
          scope. A \texttt{var} declared inside one of these blocks is
          not visible outside it.
\end{enumerate}

\subsubsection{Why Static Scoping?}

A developer can determine which variable a name refers to by reading the
source alone, without executing the program. Dynamic scoping would make
name resolution call-order--dependent. In a pipeline context this is
especially dangerous: a function called from two different
\texttt{on\_failure} handlers could see different variable bindings
depending on which handler invoked it. This makes the same source code
mean different things in different call contexts --- precisely the kind
of unreliability WorkflowScript is designed to prevent.

\subsubsection{What Would Break Under Dynamic Scoping?}

Consider a function \texttt{fun check(t: task) -> bool} that reads a
top-level variable \texttt{var threshold: int}. With static scoping,
\texttt{threshold} always resolves to the top-level declaration,
regardless of where \texttt{check} is called. With dynamic scoping, if
a caller had a local variable named \texttt{threshold}, the function
would silently read the caller's value instead. Pipeline authors would
have to track all active variable names at every call site --- an
unreasonable cognitive burden.

\subsection{Lifetime}

\begin{table}[h!]
\centering
\begin{tabular}{@{}p{4cm}p{4.5cm}p{5.5cm}@{}}
\toprule
\textbf{Entity} & \textbf{Lifetime} & \textbf{Rationale} \\
\midrule
\texttt{task} declarations
    & Static --- entire program execution
    & Stage references must remain valid whenever a pipeline runs \\
\texttt{pipeline} declarations
    & Static --- entire program execution
    & Pipelines may be invoked anywhere in the program \\
Top-level \texttt{var} declarations
    & Static --- entire program execution
    & Top-level variables function as global pipeline configuration \\
Function parameters
    & Stack-dynamic --- created on call, destroyed on return
    & Standard; no reason to outlive the invocation \\
Local \texttt{var} inside a block
    & Stack-dynamic --- created on block entry, destroyed on exit
    & Standard block-scoped allocation \\
\bottomrule
\end{tabular}
\caption{Lifetime of entities in WorkflowScript}
\end{table}

\subsubsection{Why Static Lifetime for Tasks?}

A pipeline's \texttt{stage} declarations are resolved at compile time,
but the pipeline may be executed (via \texttt{run}) at any point during
program execution. If tasks were stack-dynamic, a task declared inside a
function could be destroyed before a pipeline referencing it runs. Making
task lifetime static ensures that any \texttt{stage} reference that
passes the type checker is also safe at runtime --- the referenced task
will always exist.

\subsubsection{Why Stack-Dynamic for Local Variables?}

Local variables need to hold different values across different function
calls and loop iterations. Stack-dynamic allocation achieves this
naturally: each call to a function creates a fresh binding for each
parameter and local variable. Explicit-heap-dynamic allocation would add
unnecessary complexity (requiring manual deallocation or a garbage
collector) for values that are structurally bounded by the block they
live in.

% -----------------------------------------------------------------------
\section{Expressions \& Assignment}
% -----------------------------------------------------------------------

\subsection{Operator Precedence and Associativity}

The table below lists all binary operators from lowest to highest binding strength. All
binary operators are left-associative. Precedence is encoded directly into the
recursive-descent grammar (Section~4.3) --- each level calls the level above it --- so the
table is a reading aid, not a separate mechanism.

\begin{table}[h!]
\centering
\begin{tabular}{@{}lll@{}}
\toprule
\textbf{Precedence} & \textbf{Operators} & \textbf{Associativity} \\
\midrule
1 (lowest) & \texttt{||} & left \\
2          & \texttt{\&\&} & left \\
3          & \texttt{==} \quad \texttt{!=} & left \\
4          & \texttt{<} \quad \texttt{<=} \quad \texttt{>} \quad \texttt{>=} & left \\
5          & \texttt{+} \quad \texttt{-} & left \\
6          & \texttt{*} \quad \texttt{/} & left \\
7 (highest) & unary \texttt{-} \quad \texttt{!} & right (prefix) \\
\bottomrule
\end{tabular}
\caption{Operator precedence and associativity (lowest to highest)}
\end{table}

The unary operators bind more tightly than any binary operator. \texttt{!a \&\& b} parses
as \texttt{(!a) \&\& b}, and \texttt{-a * b} parses as \texttt{(-a) * b}.

\subsection{Short-Circuit Evaluation of Boolean Operators}

WorkflowScript uses \textbf{short-circuit (lazy) evaluation} for \texttt{\&\&} and
\texttt{||} (Sebesta \S7.6):

\begin{itemize}
    \item \texttt{e$_1$ \&\& e$_2$}: if \texttt{e$_1$} evaluates to \texttt{false},
          \texttt{e$_2$} is not evaluated and the result is \texttt{false}.
    \item \texttt{e$_1$ || e$_2$}: if \texttt{e$_1$} evaluates to \texttt{true},
          \texttt{e$_2$} is not evaluated and the result is \texttt{true}.
\end{itemize}

Sebesta \S7.6 motivates this choice with the canonical array-bounds guard example: without
short-circuit evaluation, the second operand of \texttt{\&\&} would be evaluated even when
the first is false, causing a potential runtime error. WorkflowScript adopts short-circuit
semantics to support the same guard pattern in \texttt{while} conditions.

Sebesta \S7.6 also warns that short-circuit evaluation interacts with side effects: if the
unevaluated operand contains a side-effecting expression, the side effect is silently skipped.
WorkflowScript mitigates this by restricting side effects: the only side-effecting statement
is \texttt{run}, which cannot appear inside an expression. Function calls in expressions
return values but do not modify global state, so the skip-on-short-circuit hazard does not
apply.

The type checker evaluates both operands unconditionally for type correctness. This is sound
because type checking is a static pass that does not execute expressions.

\subsection{Order of Operand Evaluation}

Sebesta \S7.2.2 identifies operand evaluation order as a design issue whenever operands can
have side effects: if a function call in an expression modifies a variable that also appears
as another operand, the result depends on which operand is evaluated first. Two solutions
exist: disallow functional side effects, or define a specific evaluation order
(Sebesta \S7.2.2.1).

WorkflowScript chooses the second option. Operand evaluation order is \textbf{defined as
left-to-right}: for a binary expression \texttt{e$_1$ op e$_2$}, \texttt{e$_1$} is always
evaluated before \texttt{e$_2$}. This matches the approach Java takes (Sebesta \S7.2.2.1:
``The Java language definition guarantees that operands appear to be evaluated in
left-to-right order'') and is consistent with short-circuit semantics, which requires the
left operand to be known before deciding whether to evaluate the right.

\subsection{Assignment: Statement or Expression?}

Assignment (\texttt{x = expr}) is a \textbf{statement}, not an expression
(Sebesta \S7.7.1, \S7.7.5). It cannot appear inside a larger expression; it must stand
alone on its own line. The parser enforces this by distinguishing \texttt{IDENT "="} (assignment
statement) from \texttt{IDENT "("} (call expression) using one token of lookahead.

Sebesta \S7.7.5 discusses the C design where assignment is an expression, and identifies its
primary hazard: \texttt{if (x = y)} is silently accepted when the programmer intended
\texttt{if (x == y)}. WorkflowScript structurally prevents this class of bug --- \texttt{=}
is never valid in an expression context, so the mistake is a parse error. Chained assignment
(\texttt{a = b = 5}) is not supported, consistent with Fortran and Ada's approach
(Sebesta \S7.7.1).

% -----------------------------------------------------------------------
\section{Design Rationale}
% -----------------------------------------------------------------------

For each of the five numbered design decisions from sections 4.5--4.7 above, the following
paragraphs justify the choice and name at least one trade-off accepted.

\subsection{Decision 1 --- Strong Static Typing}

WorkflowScript is strongly typed because reliability is its primary language criterion. A
pipeline configuration that silently accepts the wrong type of value can fail at execution
time, potentially after a long-running job has already started. By detecting all type errors
before execution begins, WorkflowScript guarantees that a program which passes the type checker
will not crash due to a type mismatch. The trade-off accepted is reduced flexibility: users
cannot, for instance, write a function that accepts ``any numeric type'' without specifying
\texttt{int} or \texttt{float}. This is acceptable because the domain is narrow --- task
parameters are always concrete integers or booleans --- and generics would significantly
complicate the type checker without enabling any new workflow patterns.

\subsection{Decision 2 --- int-to-float Widening Only}

Allowing \texttt{int} to widen to \texttt{float} eliminates friction when computing derived
metrics (e.g., \texttt{t.timeout * 1.5} for a soft deadline). Disallowing
\texttt{float}-to-\texttt{int} narrowing prevents silent truncation: \texttt{retries: 2.7}
is meaningless and should be a compile-time error. The trade-off is that converting a float
result back to an integer must be done explicitly, which is not currently supported in the
language. This is intentional: if a programmer needs \texttt{int(x / 3.0)}, they should be
forced to think about what truncation means for their pipeline logic. Keeping coercion
one-directional keeps the \texttt{typesCompatible} function simple and its behavior obvious.

\subsection{Decision 3 --- Name Equivalence for \texttt{task}}

Name equivalence was chosen because \texttt{task} is a domain role, not a data structure.
A \texttt{build} task and a \texttt{test} task carry the same three fields, but they represent
fundamentally different units of work. If structural equivalence were used, any ad-hoc record
with \texttt{timeout: int}, \texttt{retries: int}, \texttt{parallel: bool} fields would be
interchangeable with a declared task --- which would allow a programmer to accidentally pass
the wrong task to a pipeline stage. With name equivalence and a single \texttt{task} type,
all tasks are interchangeable at the type level (they all have type \texttt{task}), but their
identities are preserved at the value level (the runtime stores each task under its declared
name). The trade-off is that the language cannot distinguish \texttt{task build} from
\texttt{task deploy} in a function signature; both accept any \texttt{task}. A richer type
system with per-task singleton types would fix this but is beyond the scope of the project.

\subsection{Decision 4 --- Short-Circuit Boolean Evaluation}

Short-circuit evaluation is the natural choice for a language where the right-hand side of
\texttt{\&\&} or \texttt{||} may be an expression with a side effect or a potential runtime
error (e.g., division by zero in a guard condition). It matches the expectations of
programmers coming from Go, Python, C, and most mainstream languages. The implementation is
clean: \texttt{evalLogical} is a separate method from \texttt{evalBinary}, called specifically
when the operator is \texttt{\&\&} or \texttt{||}. The trade-off is a small asymmetry: the
type checker always evaluates both operands for type correctness, while the interpreter may
not evaluate the right operand. This means a type error in the right operand of
\texttt{false \&\& <bad\_expr>} is caught at compile time even though the expression would
never be evaluated at runtime. This is the correct behavior --- type checking and execution
are separate passes.

\subsection{Decision 5 --- Assignment as Statement, Not Expression}

Treating assignment as a statement eliminates an entire class of bugs. The most common is
\texttt{if x = 5} where the programmer intended \texttt{if x == 5}: in languages where
assignment is an expression, the if-condition is silently the assigned value (\texttt{5},
which is truthy), not the comparison result. WorkflowScript's parser rejects this pattern
structurally: \texttt{=} is never parsed inside an expression context. The trade-off is that
chained assignment (\texttt{a = b = 5}) is unsupported, but chaining is rarely needed in
workflow configuration scripts and its absence does not reduce expressiveness in the domain.

\end{document}
